% Standalone review; generated output belongs outside the repository.
% latexmk -pdf -interaction=nonstopmode -halt-on-error \
%   -outdir=/tmp/sfcc-refinement-obstruction-review \
%   docs/research/sfcc-refinement-obstruction-review.tex
\documentclass[11pt]{article}
\usepackage[T1]{fontenc}
\usepackage[utf8]{inputenc}
\usepackage{lmodern}
\usepackage[margin=1in]{geometry}
\usepackage{amsmath,amssymb,amsthm,mathtools}
\usepackage{booktabs,array,enumitem,xurl}
\usepackage[colorlinks=true,linkcolor=blue,urlcolor=blue]{hyperref}
\setlist{itemsep=3pt,topsep=5pt}
\setlength{\emergencystretch}{2em}
\newcommand{\Q}{\mathbb Q}
\newcommand{\R}{\mathbb R}
\newcommand{\conv}{\operatorname{conv}}
\title{SFCC Refinement Obstruction: Independent Review}
\author{Galacticad research documentation}
\date{October 13, 2026}
\begin{document}
\maketitle

\paragraph{Verdict and scope.}
Result B of \path{docs/research/sfcc-correctness-proofs.tex} is correct
under its stated exact-geometric premises: the specified tetrahedron has
no finite face-$2{:}1$ balanced dyadic octree whose leaves all satisfy the
one-edge-or-one-corner rule. The $24$-neighbor obstruction is also correct
for distinct straight edges with generic face crossings. Neither result
proves that the finite-budget Rust exporter runs forever, nor certifies
its output. This review does not validate Results A, C, or D or the
curved-surface estimates. The pre-existing proof source was left unchanged.
Source mapping below was inspected at repository commit
\texttt{ce1c870}; this document and the proof are separate working-tree
research artifacts.

\section{Premises required by the obstruction}

The mathematical input is a bounded, full-dimensional rational convex
polyhedron, or a regularized Boolean combination of such polyhedra, with
an embedded closed manifold boundary. Its complete geometric feature
graph consists of genuine straight crease segments and their corners.
The partition consists of axis-aligned dyadic cubes in one fixed root;
leaf interiors are disjoint. Every surface-intersecting leaf must be
accepted. Acceptance allows at most one feature edge, or exactly one
corner and only edges incident to that corner. Positive-area face
neighbors differ in depth by at most one. No unrestricted local
arrangement or other fallback may replace this acceptance requirement.

These are substantive premises. A missed feature, a deleted surface
cell, a depth cap, or a fallback allowing several nonincident curves is
not a proof of finite acceptance under the rule. Conversely, the rule
does not prohibit meshing the solid by a different construction.

\section{Independent check of the tetrahedron}

Put $v=(1/3,1/3,1/3)$ and
\[
 w_1=(0,1,-61),\quad w_2=(1,0,-61),\quad
 w_3=(-3,-2,-61),\qquad b_i=v+w_i/366.
\]
All vertices lie strictly inside $[0,1]^3$. The base height is $1/6$,
and
\[
 \det(b_1-v,b_2-v,b_3-v)=\frac{1}{133956}\ne0.
\]
The solid is a nondegenerate convex tetrahedron; its boundary is an
embedded sphere, and all six edges are genuine creases. No Boolean
regularization subtlety is needed for this witness.

\paragraph{The all-depth step.}
Depths $0$ and $1$ cannot contain an accepted apex leaf because the
apex cell also contains all three base corners. For an apex leaf at
depth $d\ge2$, let $h=2^{-d}$, $z_0=h\lfloor 2^d/3\rfloor$, and
$\Delta=1/3-z_0$. Since $2^d$ alternates between $1$ and $2$ modulo $3$,
\[
 \Delta/h\in\{1/3,2/3\}.
\]
This is an induction on a two-state recurrence, not an inference from
checking finitely many depths. Each horizontal apex coordinate is at
least $h/3$ from a side and $h/6$ from the face midline. The exit point
of edge $i$ on the lower face is
\[
 p_i=\left(\frac13+\frac{\Delta w_{ix}}{61},
            \frac13+\frac{\Delta w_{iy}}{61},z_0\right).
\]
Its maximum horizontal displacement is at most
$3\Delta/61\le2h/61<h/6$. Thus all three segments first leave the apex
cell through one and the same open quarter of the lower face.
The crossings occur before the base endpoints because
$z_0\ge1/4>1/6$.

\paragraph{Why balancing forces a contradiction.}
An adjacent leaf has side $2h$, $h$, or $h/2$. Its transverse boundary
coordinates are multiples of $h/2$; none can cut the interior of the
chosen quarter. Consequently all three exits enter a single neighboring
leaf $N$. They enter its interior because the crossings are transverse
and avoid its side boundaries. The apex is strictly above $N$.
The only endpoint common to all three apex edges is $v$, so no other
corner in $N$ can justify all three edges. The neighbor violates both
alternatives of the acceptance rule. This proves impossibility for
\emph{every finite adaptive partition}, not merely a uniform grid or
one particular refinement traversal.

\paragraph{Genericity is genuinely satisfied.}
Each $w_i$ has three pairwise distinct residues modulo $3$. If two
coordinates of $v+t w_i$ were dyadic, then
\[
 w_{ij}p_k-w_{ik}p_j=(w_{ij}-w_{ik})/3
\]
would be dyadic, although its reduced denominator is divisible by $3$.
Thus no apex edge meets any dyadic grid edge at any depth. The base
edges have nondyadic height and satisfy, respectively,
\[
 x+y=245/366,\quad x-y=-1/366,\quad x-2y=-121/366.
\]
None can have both horizontal coordinates dyadic. A parallel feature
line is disjoint from the corresponding grid planes; all its actual
face crossings are transverse. The tetrahedron obstruction does not
rely on a vertex on the grid, tangency, or numerical roundoff.

\paragraph{Quantifiers and translations.}
The explicit theorem is $\exists P$ for the fixed root $[0,1]^3$.
An axis-preserving affine map $p\mapsto o+Lp$ transfers it to any
rational root $o+[0,L]^3$: $\forall Q\ \exists P_Q$ with the same
obstruction. It does not prove that this particular tetrahedron fails
for every translated grid, nor for a sequence of grids whose origins
change with the requested maximum depth.

\section{The 24-neighbor argument}

An accepted cell containing a corner of degree at least $25$ contains
no other edge endpoint. Each incident straight segment must therefore
exit the cell. Generic crossings put each exit in a face interior.
Face-$2{:}1$ balance permits at most four neighboring leaves per face,
and hence at most $6\cdot4=24$ across the six faces. At least two
distinct edges enter the same neighbor. Distinct straight segments with
common endpoint $v$ cannot also share their other endpoint: they
would be the same geometric segment. The neighbor does not contain
$v$, so it cannot be accepted.

The straightness and genericity assumptions are necessary to this
proof as written: two curved arcs can share both endpoints, and a
grid-edge exit requires different incidence accounting. Degree $24$
is only a necessary counting bound, never a sufficiency result; the
degree-three tetrahedron already refutes sufficiency. Unlike the
tetrahedron's particular quarter argument, this counting obstruction
applies to \emph{any} grid satisfying the generic-crossing premise.

The checker includes an explicit rational witness for existence of
such corners. With $K=6\cdot997\cdot16$ and $i=-12,\ldots,12$, take
\[
 b_i=v+(3i,3i^2+1,-997)/K.
\]
These $25$ points form a strictly convex polygon in $z=31/96$;
its convex hull with $v$ is a manifold pyramid with $25$ apex edges.
The points lie inside the root. Apex directions again have distinct
residues modulo $3$. Every base edge joining indices $i,j$ satisfies
\[
 y-(i+j)x=(1-i-j)/3+(1-3ij)/K,
\]
with a nondyadic right-hand side. Together with the nondyadic base
height, this excludes all dyadic grid-edge contacts. This finite exact
witness verifies that the high-valence hypothesis is nonvacuous.

\section{What transfers to the current Rust implementation}

The following mapping uses function bodies, not documentation claims.
Paths are relative to \path{gcad-wasm/kernel/src/}.

\begin{description}[style=nextline,leftmargin=0pt,labelwidth=0pt,labelsep=0pt]
\item[Feature rule: \path{sfcc/refine_criteria.rs}, \texttt{classify\_cell\_features}.]
Two in-cell corners require a split. A corner accepts only incident
curves with one boundary crossing. A cornerless cell accepts at most
one curve with two crossings; per-face multiplicity, small tangent
dot products, and pin visibility can require further splitting.
The obstruction applies to this idealized rule if the complete graph
is discovered and the geometric queries are correct. Actual queries
are numerical and operate on compiled candidates. Acceptance caused
by omission or tolerance merging does not evade the exact theorem;
it violates a premise.

\item[Balance: \path{sfcc/octree.rs}, \texttt{ripple\_balance}.]
Face balancing is unconditional. The \texttt{enforce\_edge\_balance}
switch adds edge-neighbor balancing; switching it off does not remove
the obstruction's face-balance premise.

\item[Finite budgets: \path{sfcc/octree.rs}, \texttt{apply\_decision}.]
A requested split at \texttt{depth\_max} marks the cell degenerate and
returns. The default maximum depth is $8$, and
\path{tuning.rs} caps the lattice at $14$.
The obstruction concerns reaching an all-accepted partition, not
termination of this bounded procedure. The pipeline includes
degenerate leaves in \texttt{unresolved\_cells}; validation cannot
pass while that count is positive.

\item[Grid: \path{sfcc/pipeline.rs}, \texttt{build\_pipeline\_context}.]
The root is padded, then shifted on axis $i$ by
$j_i=c_i L2^{-D}/4$, with
$c_i\in\{\sqrt2-1,\sqrt3-1,\sqrt5-2\}$ and
$D=\texttt{depth\_max}$. A run uses one fixed hierarchy, but changing
$D$ changes its origin. The exact fixed-root example is therefore
not an all-budget floating-point runtime fixture. The implemented
values are f64 approximations, not irrational numbers. Generic grid
placement alone does not solve the high-valence counting obstruction.

\item[Corner exemption: \path{sfcc/pipeline.rs}, \texttt{decide\_cell}.]
After feature classification succeeds, a corner cell bypasses both
the sign-change gate and the smooth criteria. Consequently the
proof's additional observation about a closed triangular face trace
does not mean this implementation rejects every such apex cell by
its sign pattern. Its neighbor can still violate the feature rule.

\item[Fallback: \path{sfcc/cell_mesh.rs}, \texttt{mesh\_cells\_subset}.]
Multiple corners, certain degenerate multi-curve cells, and curved
corner arcs trigger \path{mesh_feature_graph}. This is a richer
filling rule; it can produce geometry in an unresolved leaf. It does
not clear the octree's degenerate tag. The graph routine also requires
exactly one boundary loop and is not the unrestricted arrangement
completion proposed in Result C. No success theorem for it follows
from either obstruction.

\item[Input class: \path{scene_bridge.rs}, \texttt{BridgeNode}.]
The bridge supports boxes, spheres, cylinders, cones, profile
extrusions, lofts, lathes, similarities, and Boolean/blend nodes. It
has no arbitrary rational-halfspace polyhedron or tetrahedron leaf.
In ideal real arithmetic, finitely many sufficiently large rotated
boxes can realize the tetrahedron's supporting halfspaces near the
root, hence their intersection can realize the tetrahedron. That
representation uses numerical rotations in the actual bridge and
does not preserve these exact rational coordinates automatically.
The check supplied here tests the mathematical input class; it is
not evidence that a particular JSON scene compiles its exact graph.
\end{description}

\paragraph{Angular rejection is a separate obstruction.}
For a fixed straight edge direction $t$ crossing planes normal to
$e_i$, the value $|t\cdot e_i|$ does not increase under subdivision.
Thus a nonzero crossing rejected by a fixed positive
\texttt{tangential\_epsilon} remains rejected on that family of
faces. Exact transversality is not a uniform angle bound. This is
separate from the balanced-neighbor argument; the tetrahedron's
lower-face exits are almost perpendicular to that face.

\section{Reproduction and interpretation of a passing check}

Run from the repository root:
\begin{verbatim}
python3 scripts/check-sfcc-refinement-obstruction.py
\end{verbatim}
It uses only integer arithmetic and \texttt{fractions.Fraction}.
It checks the tetrahedron volume, root inclusion, base planes,
modulo-$3$ genericity invariants, the two-state recurrence, all exits
and possible neighboring grid levels at depths $2$ through $199$,
and that every apex-cell vertex is outside. It also checks strict
convexity and genericity of the explicit $25$-gon pyramid.
The $198$ depth instances check transcription. The recurrence,
drift inequality, dyadic alignment, and pigeonhole arguments above
provide the all-depth reasoning; a finite sweep alone would not.

A successful exit means \emph{the obstruction's arithmetic was
verified}, not that the exporter succeeds. Future runtime regressions
must record returned geometry, unresolved counts, and recovery mode
separately. A still-failing correctness test should retain its desired
correctness assertion and an explicit known-failure/ignored marker
with its counterexample and reason, rather than invert the assertion
to accept wrong geometry silently. A test requiring unresolved status
is appropriate when it explicitly checks honest budget reporting.
An impossible all-accepted-partition requirement cannot become a
passing regression until at least one of the obstruction's premises
is deliberately changed.

\paragraph{Next theoretical decision.}
Keep the obstruction as a design constraint. A completion must relax
the one-feature restriction, face balance at protected cells, the
fixed-grid partition model, or the requirement that refinement alone
make every leaf simple. Merely increasing depth, improving numerical
precision, or adding a genericity shift is not a universal termination
argument for the stated rule.

\end{document}
